\documentclass[letterpaper,10pt]{article}

\usepackage[left=0.5in,top=0.5in,right=0.5in,bottom=0.5in]{geometry}
\usepackage{hyperref}
\usepackage{enumitem}
\usepackage{titlesec}
\usepackage{array}
\usepackage{fontenc}
\usepackage[utf8]{inputenc}

\hypersetup{colorlinks=false}

% Remove page numbers
\pagestyle{empty}

% Section formatting
\titleformat{\section}{\large\bfseries}{}{0em}{}[\titlerule]
\titlespacing{\section}{0pt}{6pt}{4pt}

% No paragraph indent
\setlength{\parindent}{0pt}

% Tight lists
\setlist[itemize]{noitemsep, topsep=2pt, parsep=0pt, leftmargin=1.2em}

\begin{document}

%--- HEADER ------------------------------------------------------------------
\begin{center}
  {\LARGE \textbf{Vansh Nawander}} \\[4pt]
  \href{mailto:vansh.nawander@research.iiit.ac.in}{vansh.nawander@research.iiit.ac.in} \;$|$\;
  \href{https://www.linkedin.com/in/vanshnawander/}{LinkedIn} \;$|$\;
  \href{https://github.com/vanshnawander}{GitHub} \;$|$\;
  \href{https://vanshnawander.github.io/vansh/index.html}{Website} \;$|$\;
  +91-8125544202
\end{center}

%--- EDUCATION ---------------------------------------------------------------
\section{Education}

\textbf{International Institute of Information Technology, Hyderabad (IIIT-H)} \hfill Hyderabad, India \\
\textit{M.S. by Research, Computer Science} \hfill 2025 -- Present \\
Advisor: Dr.\ Manish Shrivastava (LTRC) \quad $|$ \quad Focus: Diffusion Models, Quantization, LLM Inference Systems

\medskip
\textbf{Vardhaman College of Engineering} \hfill Hyderabad, India \\
\textit{B.Tech, Computer Science and Engineering} \hfill 2020 -- 2024

%--- RESEARCH INTERESTS ------------------------------------------------------
\section{Research Interests}

GPU architectures \& kernel optimization \textbullet{} LLM inference \& training systems \textbullet{}
Model quantization \textbullet{} Deep learning \textbullet{} Diffusion models \textbullet{} World models

%--- EXPERIENCE --------------------------------------------------------------
\section{Experience}

\textbf{AI/ML Consultant -- Alonzo AI} \hfill April 2025 -- March 2026
\begin{itemize}
  \item Fine-tuned and deployed production ASR and TTS systems for speech applications.
  \item Built an end-to-end agentic speech system integrating speech recognition, language understanding, and speech synthesis into a unified pipeline.
  \item Handled backend engineering and NLP-based search infrastructure.
\end{itemize}

\textbf{AI/ML Engineer -- Datazoic} \hfill Oct 2024 -- Jun 2025
\begin{itemize}
  \item Fine-tuned large language models on A100 GPU servers and deployed them to production.
  \item Built RAG-based question answering systems for retrieval over thousands of documents.
  \item Engineered AI agents with LangGraph for a CRM platform, enabling automated actions and improved UX.
\end{itemize}


%--- PROJECTS ----------------------------------------------------------------
\section{Projects}

\textbf{FlashInfer Kernel Contest @ MLSys 2026} \hfill
\href{https://github.com/vanshnawander/flashinfer-mlsys}{\texttt{[GitHub]}}
\begin{itemize}
  \item Competing in the FlashInfer AI Kernel Generation Contest at MLSys 2026, writing high-performance GPU kernels (fused MoE) for NVIDIA Blackwell (B200) GPUs.
  \item Implemented kernels in Triton and CUDA; benchmarked on B200 hardware via Modal cloud infrastructure. Produced a technical report analysing kernel performance against the FlashInfer-Bench evaluation suite.
\end{itemize}

\textbf{CUDA \& GPU Kernel Sessions} \hfill
\href{https://github.com/vanshnawander/cuda-sessions}{\texttt{[GitHub]}}
\begin{itemize}
  \item Implemented GPU kernels from scratch: tiled matrix multiplication, memory-optimised CUDA kernels, and Triton kernels; includes Mojo GPU puzzle solutions and GPUMode practice exercises.
  \item Stack: C++, CUDA, Triton, Mojo, Python.
\end{itemize}

\textbf{Language Model from Scratch}
\begin{itemize}
  \item Built a transformer-based language model from scratch in PyTorch, implementing multi-head attention, positional encoding, and the full training pipeline.
  \item Covered tokenisation, pretraining objectives, and scaling considerations; used as a deep-dive into LLM internals and architecture design.
\end{itemize}


%--- PUBLICATIONS ------------------------------------------------------------
\section{Publications}

\begin{itemize}
  \item V.\ Nawander et al., ``Span Identification and Summarization,''
        \textit{Workshop Proceedings, NAACL 2025}.
        \href{https://aclanthology.org/2025.cl4health-1.33.pdf}{[Link]}

\end{itemize}

%--- SKILLS ------------------------------------------------------------------
\section{Skills}

\begin{tabular}{@{} >{\bfseries}l @{\hspace{3ex}} l}
Systems \& GPU  & CUDA, Triton, C, C++, GPU kernel optimisation \\
ML \& DL        & PyTorch, JAX, vLLM, HuggingFace, Scikit-learn, NumPy, Pandas \\
LLM / AI        & LangChain, LangGraph, DSPy, ADK, RAG, fine-tuning (A100) \\
Languages       & Python, C++, C, JavaScript \\
Web / Backend   & FastAPI, Flask, React.js, Node.js, Express.js \\
Infra / DB      & Linux, Redis, Celery, AWS, MySQL, PostgreSQL, MongoDB \\
\end{tabular}

%--- ACHIEVEMENTS ------------------------------------------------------------
\section{Achievements \& Leadership}

\begin{itemize}
  \item \textbf{Smart India Hackathon 2023 -- National Winner:} Won the national-level competition for problem statement \#1384 by the Ministry of Commerce and Industry.
  \item \textbf{IndoML Datathon (NIQ Sponsored):} Won 1st and 5th prizes at the IndoML conference datathon challenge.
  \item \textbf{Microsoft ImagineCup 2024:} Qualified for the first round.
\end{itemize}

\end{document}